% ============================================================================
% 03_optimization / 04_distillation
% Section-body file: included immediately after the section heading in the
% master document.  It intentionally contains no document-level commands,
% package imports, chapter heading, section heading, or bibliography command.
% ============================================================================

The preceding discussion of optimization established how a language model
learns a predictive function by adjusting its parameters against a training
objective.  That formulation, however, leaves an important deployment
question unresolved.  A model may achieve strong predictive performance and
yet remain too expensive to store, execute, or serve at scale.  Increasing
model capacity can improve the quality of the learned distribution, but it
also increases parameter memory, arithmetic operations, inference latency,
and the resources required to serve long sequences.  The optimization problem
therefore has a second dimension: how can the information acquired by a
computationally expensive model be transferred to a model whose execution is
less demanding?

Knowledge distillation addresses this question by training a \emph{student}
model under the guidance of a \emph{teacher} model.  The teacher is usually
larger or more capable, while the student is designed to have lower storage
and inference costs.  The objective is not necessarily to copy the teacher's
parameters or architecture.  Instead, the student is trained to reproduce
selected aspects of the teacher's behavior, such as its output probability
distributions, intermediate representations, generated responses, or
reasoning traces.  The information made available by the teacher and the
student's ability to absorb it determine the form of the distillation
objective.  Knowledge distillation is consequently a family of optimization
procedures rather than a single loss function.  This characterization is
consistent with the LLM distillation surveys
\parencite{Yang2024SurveyKD,Fang2026KnowledgeDatasetDistillation}.

The distinction between model compression and knowledge transfer is central
to this chapter.  A smaller parameter count does not by itself guarantee that
the student retains the teacher's useful behavior.  The student must learn a
suitable approximation from finite data and within a finite optimization
budget.  The following development begins with the probability distributions
produced by autoregressive language models, establishes the mathematical
relationship between those distributions, and then examines the different
interfaces through which a teacher can supervise a student.

\subsection{Predictive distributions and the mathematical objective of distillation}
\label{subsec:distillation-probabilistic-foundation}

An autoregressive language model predicts a distribution over the next token
conditioned on the tokens already observed.  Let the vocabulary be
\(\mathcal V=\{v_1,\ldots,v_C\}\), where \(C\) is the vocabulary
size.  At position \(t\), the context is
\(c_t=x_{<t}=(x_1,\ldots,x_{t-1})\).  A model with parameters
\(\theta\) maps this context to a logit vector
\begin{equation}
    \mathbf z_\theta(c_t)
    =\bigl(z_{\theta,1}(c_t),\ldots,z_{\theta,C}(c_t)\bigr)^{\mathsf T}
    \in\mathbb R^C.
    \label{eq:distillation-logit-vector}
\end{equation}
The logits are unrestricted real-valued scores; they are not probabilities
and need not sum to one.  The softmax transformation converts them into a
categorical probability distribution:
\begin{equation}
    p_\theta(v_i\mid c_t)
    =\frac{\exp(z_{\theta,i}(c_t))}
    {\displaystyle\sum_{j=1}^{C}\exp(z_{\theta,j}(c_t))}.
    \label{eq:distillation-softmax}
\end{equation}
The resulting probability vector belongs to the probability simplex
\begin{equation}
    \Delta^{C-1}
    =\left\{\mathbf p\in\mathbb R^C:
    p_i\geq 0,\;\sum_{i=1}^{C}p_i=1\right\}.
    \label{eq:distillation-simplex}
\end{equation}

For a sequence \(\mathbf x=(x_1,\ldots,x_T)\), the autoregressive
factorization gives
\begin{equation}
    P_\theta(\mathbf x)
    =\prod_{t=1}^{T}P_\theta(x_t\mid x_{<t}).
    \label{eq:distillation-autoregressive-factorization}
\end{equation}
Taking the negative logarithm converts the product into a sum.  The
next-token negative log-likelihood (NLL) objective is therefore
\begin{equation}
    \mathcal L_{\mathrm{LM}}(\theta)
    =-\frac{1}{T}\sum_{t=1}^{T}
    \log P_\theta(x_t\mid x_{<t}).
    \label{eq:distillation-nll}
\end{equation}
This is also the token-averaged cross-entropy between the empirical target
and the model prediction.  To see this, let \(\mathbf e_{x_t}\) be the
one-hot vector whose component corresponding to the observed token is one.
Then
\begin{equation}
    \mathcal L_{\mathrm{NLL}}^{(t)}
    =-\sum_{i=1}^{C}e_{x_t,i}
    \log p_\theta(v_i\mid c_t)
    =-\log p_\theta(x_t\mid c_t).
    \label{eq:distillation-one-hot-ce}
\end{equation}
A one-hot target identifies the observed token, but it does not directly
express how plausible the model considers the other vocabulary items.  A
teacher's full predictive distribution contains this additional information.
For example, suppose the observed token is \(v_1\) and the teacher's
probability vector is
\begin{equation}
    \mathbf p_T=(0.70,0.20,0.08,0.02)^{\mathsf T}.
    \label{eq:distillation-dark-knowledge-example}
\end{equation}
The distribution communicates not only that \(v_1\) is preferred, but
also that \(v_2\) is considered more plausible than \(v_3\), and
\(v_3\) more plausible than \(v_4\).  These relative preferences
among alternatives are commonly called the teacher's
\emph{dark knowledge} in the distillation literature
\parencite{Yang2024SurveyKD,Fang2026KnowledgeDatasetDistillation}.

The uncertainty represented by a categorical distribution can be quantified
using Shannon entropy.  For a probability vector \(\mathbf p\),
\begin{equation}
    H(\mathbf p)=-\sum_{i=1}^{C}p_i\log p_i.
    \label{eq:distillation-entropy}
\end{equation}
A distribution concentrated on one category has low entropy, whereas the
uniform distribution has entropy \(\log C\).  Thus, teacher predictions
with the same most probable token can still convey different information
when their probability mass is distributed differently across the remaining
vocabulary.  Distillation uses this richer distribution as a learning target
rather than reducing teacher behavior to its top prediction.

Let \(\mathbf p_T(\cdot\mid c)\) and
\(\mathbf p_S(\cdot\mid c)\) denote teacher and student distributions
at the same context.  Their cross-entropy is
\begin{equation}
    H(\mathbf p_T,\mathbf p_S)
    =-\sum_{i=1}^{C}p_{T,i}\log p_{S,i},
    \label{eq:distillation-cross-entropy}
\end{equation}
whereas their Kullback--Leibler (KL) divergence is
\begin{equation}
    D_{\mathrm{KL}}(\mathbf p_T\Vert\mathbf p_S)
    =\sum_{i=1}^{C}p_{T,i}
    \log\frac{p_{T,i}}{p_{S,i}}.
    \label{eq:distillation-forward-kl}
\end{equation}
Expanding the logarithm gives
\begin{align}
    D_{\mathrm{KL}}(\mathbf p_T\Vert\mathbf p_S)
    &=\sum_i p_{T,i}\log p_{T,i}
      -\sum_i p_{T,i}\log p_{S,i}\nonumber\\
    &=-H(\mathbf p_T)+H(\mathbf p_T,\mathbf p_S).
    \label{eq:distillation-kl-cross-entropy-identity}
\end{align}
Consequently,
\begin{equation}
    H(\mathbf p_T,\mathbf p_S)
    =H(\mathbf p_T)+D_{\mathrm{KL}}(\mathbf p_T\Vert\mathbf p_S).
    \label{eq:distillation-ce-kl-identity}
\end{equation}
When the teacher is fixed, \(H(\mathbf p_T)\) is constant with respect
to the student parameters.  Minimizing cross-entropy against the teacher is
therefore equivalent to minimizing forward KL divergence.  KL divergence is
nonnegative and is zero if the two distributions agree, subject to the usual
support conditions.  It is not symmetric: in general,
\(D_{\mathrm{KL}}(P\Vert Q)\neq D_{\mathrm{KL}}(Q\Vert P)\).

This asymmetry changes the approximation favored when the student cannot
represent the teacher distribution exactly.  Forward KL,
\(D_{\mathrm{KL}}(\mathbf p_T\Vert\mathbf p_S)\), weights the
alternatives according to teacher probability and penalizes the student for
failing to cover teacher-supported alternatives.  Reverse KL,
\(D_{\mathrm{KL}}(\mathbf p_S\Vert\mathbf p_T)\), weights the
alternatives according to student probability and penalizes student mass on
alternatives that the teacher considers unlikely.  In a multimodal target and
a restricted student family, the two objectives can favor different
approximations.  The Jensen--Shannon divergence provides a symmetric
alternative.  With
\(\mathbf m=(\mathbf p_T+\mathbf p_S)/2\), it is defined as
\begin{equation}
    D_{\mathrm{JS}}(\mathbf p_T\Vert\mathbf p_S)
    =\frac{1}{2}D_{\mathrm{KL}}(\mathbf p_T\Vert\mathbf m)
    +\frac{1}{2}D_{\mathrm{KL}}(\mathbf p_S\Vert\mathbf m).
    \label{eq:distillation-js}
\end{equation}
The divergence is symmetric and bounded.  The choice of discrepancy is thus
part of the optimization design, not a notational variation
\parencite{Song2026OnPolicyDistillation}.

\paragraph{Temperature-scaled targets.}
The teacher's original distribution may be sharply concentrated.  Temperature
scaling changes its concentration by dividing logits by a positive parameter
\(\tau\) before applying softmax:
\begin{equation}
    p_{T,i}^{(\tau)}
    =\frac{\exp(z_{T,i}/\tau)}
    {\displaystyle\sum_{j=1}^{C}\exp(z_{T,j}/\tau)},
    \qquad
    p_{S,i}^{(\tau)}
    =\frac{\exp(z_{S,i}/\tau)}
    {\displaystyle\sum_{j=1}^{C}\exp(z_{S,j}/\tau)}.
    \label{eq:distillation-temperature-softmax}
\end{equation}
For \(0<\tau<1\), differences between logits are amplified relative to
the temperature and the distribution becomes sharper.  For \(\tau>1\),
the distribution becomes smoother.  At \(\tau=1\), the original softmax
is recovered.  If the maximum logit is unique,
\begin{equation}
    \lim_{\tau\to 0^+}p_{T,i}^{(\tau)}
    =\begin{cases}
        1,&i=\arg\max_j z_{T,j},\\
        0,&\text{otherwise},
    \end{cases}
    \qquad
    \lim_{\tau\to\infty}p_{T,i}^{(\tau)}=\frac{1}{C}.
    \label{eq:distillation-temperature-limits}
\end{equation}
Temperature therefore controls how much of the teacher's relative preference
among nonmaximal alternatives is exposed to the student.

For the temperature-scaled cross-entropy
\begin{equation}
    \mathcal L_{\mathrm{KD}}^{(\tau)}
    =-\sum_{i=1}^{C}p_{T,i}^{(\tau)}
    \log p_{S,i}^{(\tau)},
    \label{eq:distillation-temperature-ce}
\end{equation}
differentiation with respect to a student logit gives
\begin{equation}
    \frac{\partial\mathcal L_{\mathrm{KD}}^{(\tau)}}
    {\partial z_{S,j}}
    =\frac{1}{\tau}
    \left(p_{S,j}^{(\tau)}-p_{T,j}^{(\tau)}\right).
    \label{eq:distillation-temperature-gradient}
\end{equation}
In the high-temperature regime, the difference between the softened
probabilities is itself approximately proportional to \(1/\tau\).  The
unscaled gradient can therefore shrink approximately as \(1/\tau^2\).
A conventional formulation compensates for this scale change by multiplying
the KL term by \(\tau^2\):
\begin{equation}
    \mathcal L_{\mathrm{KD}}
    =\tau^2D_{\mathrm{KL}}
    \left(\mathbf p_T^{(\tau)}\Vert
    \mathbf p_S^{(\tau)}\right).
    \label{eq:distillation-temperature-scaled-kl}
\end{equation}
The factor is a gradient-scale compensation, not a requirement that every
distillation implementation must use the same temperature convention.  The
temperature also controls the dispersion of the transferred distribution and
is discussed in Bayesian interpretations of distillation
\parencite{Fang2026KnowledgeDatasetDistillation,Yang2024SurveyKD}.

\paragraph{The optimization objective.}
Let \(f_T(c;\theta_T)\) and \(f_S(c;\theta_S)\) denote teacher and
student logits.  In conventional distillation, the teacher parameters
\(\theta_T\) are fixed and the student parameters \(\theta_S\) are
optimized.  For a context distribution \(\mathcal D\), the basic
logits-based objective is
\begin{equation}
    \theta_S^*
    =\arg\min_{\theta_S}
    \mathbb E_{c\sim\mathcal D}
    \left[
      D_{\mathrm{KL}}\left(
      \mathbf p_T^{(\tau)}(\cdot\mid c)
      \Vert
      \mathbf p_S^{(\tau)}(\cdot\mid c;\theta_S)
      \right)
    \right].
    \label{eq:distillation-basic-objective}
\end{equation}
The teacher provides a target distribution; the student receives gradients
through its own predictions.  Distillation does not replace gradient-based
optimization.  It changes the information that defines the target of that
optimization.

The original next-token target can be retained because a teacher may be
incorrect, biased, or poorly calibrated.  A common combined objective is
\begin{equation}
    \mathcal L_{\mathrm{total}}
    =(1-\alpha)\mathcal L_{\mathrm{LM}}
    +\alpha\mathcal L_{\mathrm{KD}},
    \qquad 0\leq\alpha\leq 1,
    \label{eq:distillation-combined-objective}
\end{equation}
where \(\alpha\) controls the relative contribution of the original
language-modeling target and teacher supervision.  At \(\alpha=0\), the
student is trained only on the observed tokens; at \(\alpha=1\), the
specified objective uses only the distillation term.  Intermediate values
combine both sources of supervision.  This formulation is the basis for the
pre-training distillation experiments discussed later
\parencite{Peng2024PretrainingDistillation}.

\subsection{Knowledge interfaces: white-box and black-box distillation}
\label{subsec:distillation-knowledge-interfaces}

The output distribution is only one possible representation of a teacher's
knowledge.  A language model's predictive behavior is produced by a sequence
of internal transformations, and information relevant to the final prediction
is distributed across its hidden states, attention computations, and output
layer.  The form of distillation is therefore constrained by which parts of
the teacher are accessible.  The survey by \textcite{Yang2024SurveyKD}
distinguishes white-box methods, which can access internal model information,
from black-box methods, which rely on observable outputs such as generated
text.

\paragraph{Logits-based white-box distillation.}
Logits-based distillation transfers the teacher's final vocabulary
probabilities or logits.  At each token position, the student is trained to
minimize a discrepancy such as the forward KL divergence in
\cref{eq:distillation-basic-objective}.  It provides information about the
teacher's relative preferences over alternative tokens without requiring
the student to reproduce the teacher's hidden-state dimensions or layer
structure.  Its limitation is that supervision is applied at the output
interface.  Different internal computations can produce similar output
distributions, so output matching alone does not force the student to
reproduce the teacher's intermediate representations.

\paragraph{Hint-based and representation distillation.}
Hint-based methods add supervision at intermediate layers.  Let the teacher
hidden representation at a selected layer be
\(H_T^{(l)}\in\mathbb R^{T\times d_T}\), and let the corresponding
student representation be
\(H_S^{(k)}\in\mathbb R^{T\times d_S}\).  The sequence length is
\(T\), while \(d_T\) and \(d_S\) are the teacher and student hidden
widths.  Since the dimensions and layer counts need not agree, a projection
\(\phi:\mathbb R^{d_S}\to\mathbb R^{d_T}\) can map the student
features into the teacher feature space.  A squared-error feature objective
is
\begin{equation}
    \mathcal L_{\mathrm{feat}}
    =\frac{1}{T d_T}
    \left\|H_T^{(l)}-\phi(H_S^{(k)})\right\|_F^2,
    \label{eq:distillation-feature-loss}
\end{equation}
where the squared Frobenius norm is the sum of squared entries of a matrix.  This
objective aligns selected representations, but it also imposes an
architectural correspondence that must be chosen carefully.  Matching
identically indexed layers is not automatically appropriate when teacher
and student depths differ.

Patient knowledge distillation uses several teacher layers rather than
only the final layer.  TinyBERT aligns several
levels of the Transformer computation, including embedding outputs, hidden
states, and attention-related representations.  MiniLM focuses on relations
encoded by self-attention computations, allowing transfer through attention
relationships rather than requiring every hidden coordinate to match.  These
methods illustrate that representation transfer can be designed around the
structures that are most useful to preserve, but retaining intermediate
activations increases memory and computation during teacher inference
\parencite{Yang2024SurveyKD}.

\paragraph{Black-box distillation through generated text.}
A black-box teacher exposes responses but not necessarily its logits,
parameters, or intermediate activations.  The student must then learn from
examples generated by the teacher.  In-context learning distillation uses a
task instruction and a set of demonstrations.  Let
\(D_K=\{(x_i,y_i)\}_{i=1}^{K}\) be the demonstrations and let
\(I\) be the instruction.  The teacher generates a response according to
\begin{equation}
    \widehat y
    =T\left(I,(x_1,y_1),\ldots,(x_K,y_K),x_{K+1}\right).
    \label{eq:distillation-icl-generation}
\end{equation}
The generated input--output examples are subsequently used to fine-tune the
student.  The student is thereby trained on demonstrations of how the teacher
responds to instructions and contextual examples, without requiring access
to the teacher's internal probability distribution.

Chain-of-thought distillation extends the target from a final answer to an
intermediate rationale and a final answer.  Let \(r_i\) be a teacher-generated
rationale for input \(x_i\), and let \(y_i\) be its final answer.  The student
can be trained to model the joint conditional sequence:
\begin{equation}
    \mathcal L_{\mathrm{CoT}}
    =-\frac{1}{N}\sum_{i=1}^{N}
    \log P_{\theta_S}(r_i,y_i\mid x_i).
    \label{eq:distillation-cot-loss}
\end{equation}
For an autoregressive model, the joint probability factorizes into the
conditional probabilities of the rationale tokens followed by the answer
tokens.  Rationale supervision therefore teaches an intermediate sequence,
not merely the final output.  Its effectiveness depends on the correctness
and usefulness of the teacher's rationales; a fluent explanation is not by
itself evidence that the reasoning is valid
\parencite{Yang2024SurveyKD,Fang2026KnowledgeDatasetDistillation}.

Instruction-following distillation constructs examples containing an
instruction \(I_i\), an optional input \(x_i\), and a teacher response
\(y_i^T\).  The student objective is
\begin{equation}
    \mathcal L_{\mathrm{inst}}
    =-\frac{1}{N}\sum_{i=1}^{N}
    \log P_{\theta_S}(y_i^T\mid I_i,x_i).
    \label{eq:distillation-instruction-loss}
\end{equation}
The diversity and difficulty of the instruction set determine the range of
behaviors represented in the training data.  A narrow collection of simple
instructions may yield a student that reproduces those examples but does not
generalize to complex or unfamiliar requests.  Self-instruction approaches
expand the dataset by asking a teacher to generate additional instructions
and corresponding responses, followed by quality filtering
\parencite{Yang2024SurveyKD}.

\paragraph{Uncertainty-aware and multi-teacher distillation.}
A single predictive distribution does not necessarily express all forms of
uncertainty relevant to a task.  Uncertainty-aware distillation can use
ensembles or probabilistic teachers to expose predictive variability and
train the student to approximate a teacher predictive distribution.  At a
state \(x\), a generic objective remains
\begin{equation}
    \mathcal L_{\mathrm{unc}}
    =D_{\mathrm{KL}}\left(
    p_T(\cdot\mid x)\Vert p_S(\cdot\mid x)
    \right),
    \label{eq:distillation-uncertainty-loss}
\end{equation}
but the teacher distribution can now reflect uncertainty represented by
multiple models or posterior samples.  The Bayesian interpretation discussed
in the survey treats teacher-induced regularization as information about a
prior over student behavior; temperature controls the dispersion of the
teacher-informed target.  The additional sampling or ensemble computation
must be weighed against the value of improved calibration and robustness
\parencite{Fang2026KnowledgeDatasetDistillation}.

For \(K\) teachers with predictive distributions \(p_k(y\mid x)\), a
simple multi-teacher target is the weighted mixture
\begin{equation}
    \bar p(y\mid x)
    =\sum_{k=1}^{K}w_k(x)p_k(y\mid x),
    \qquad w_k(x)\geq 0,\quad\sum_{k=1}^{K}w_k(x)=1.
    \label{eq:distillation-multi-teacher-mixture}
\end{equation}
The student can be trained to match \(\bar p\).  Uniform weights assume
equal teacher usefulness, whereas input-dependent weights can reflect
specialization or reliability.  Multi-teacher distillation can combine
complementary expertise, but it increases teacher inference cost and can
introduce conflicts among teachers with different domains, vocabularies, or
prediction preferences.  Dynamic and adaptive methods allow the supervisor
or the relative contributions of teachers to change during training.  Such
methods can support evolving data or continual learning, but co-evolution
can make convergence and error control more difficult
\parencite{Fang2026KnowledgeDatasetDistillation}.

\begin{figure}[htbp]
    \centering
    % Output-level teacher--student distillation diagram.
% Uses the standard LaTeX picture environment; no additional package is needed.
\begingroup
\setlength{\unitlength}{0.75mm}
\begin{picture}(150,62)
    \put(1,39){\framebox(25,12){\shortstack{Input\\context}}}
    \put(37,47){\framebox(30,12){Teacher}}
    \put(37,20){\framebox(30,12){Student}}
    \put(80,47){\framebox(28,12){\shortstack{Teacher\\distribution}}}
    \put(80,20){\framebox(28,12){\shortstack{Student\\distribution}}}
    \put(120,33){\framebox(27,12){\shortstack{Distillation\\objective}}}
    \put(26,45){\vector(1,1){10}}
    \put(26,45){\vector(1,-2){10}}
    \put(67,53){\vector(1,0){13}}
    \put(67,26){\vector(1,0){13}}
    \put(108,53){\vector(1,-1){11}}
    \put(108,26){\vector(1,1){11}}
    \put(2,8){\makebox(105,5)[l]{\scriptsize The teacher and student receive the same context.}}
    \put(2,2){\makebox(145,5)[l]{\scriptsize The loss compares predictive distributions; student parameters receive the update.}}
\end{picture}
\endgroup

    \caption{Output-level knowledge distillation.  The same context is
    presented to teacher and student models.  Their predictive distributions
    are compared by a distillation objective, which supplies the student
    update without requiring the student to copy the teacher parameters.}
    \label{fig:distillation-teacher-student}
\end{figure}

\subsection{Post-training and pre-training distillation}
\label{subsec:distillation-training-stages}

Distillation can be applied at different stages of language-model
construction.  Post-training distillation generally starts from a pretrained
student and transfers task-specific behavior through teacher-generated
instructions, answers, or rationales.  Pre-training distillation instead
uses teacher predictions over the pre-training corpus, so that the student
receives a teacher-derived learning signal while acquiring general language
modeling capabilities.  The two settings differ in their data distributions,
objectives, target-generation costs, and intended transfer
\parencite{Peng2024PretrainingDistillation,Yang2024SurveyKD}.

In post-training distillation, let
\(D_T=\{(x_i,y_i^T)\}_{i=1}^{N}\) be a collection of prompts and
teacher-generated responses.  Teacher-forced student training minimizes
\begin{equation}
    \mathcal L_{\mathrm{post}}
    =-\frac{1}{N}\sum_{i=1}^{N}
    \log P_{\theta_S}(y_i^T\mid x_i).
    \label{eq:distillation-posttraining-loss}
\end{equation}
For a response \(y_i^T=(y_{i,1}^T,\ldots,y_{i,L_i}^T)\), the
conditional sequence probability is
\begin{equation}
    P_{\theta_S}(y_i^T\mid x_i)
    =\prod_{t=1}^{L_i}
    P_{\theta_S}(y_{i,t}^T\mid x_i,y_{i,<t}^T).
    \label{eq:distillation-response-factorization}
\end{equation}
This procedure is applicable when the teacher is accessible only through a
text-generation interface.  Its limitation is that the student is trained on
fixed response prefixes.  During deployment, the student conditions on its
own previous predictions, which can take it away from the histories present
in the teacher demonstrations.  This discrepancy motivates on-policy
distillation, developed in \cref{subsec:distillation-on-policy}.

In pre-training distillation, teacher and student process contexts from a
pre-training corpus.  Let \(x_{<t}\) be a corpus prefix and let
\(p_T(\cdot\mid x_{<t})\) and \(p_S(\cdot\mid x_{<t})\) be their
next-token distributions.  A generic objective combines the original token
loss with a teacher-matching term:
\begin{equation}
    \mathcal L_{\mathrm{PD}}
    =(1-\alpha)\mathcal L_{\mathrm{LM}}
    +\alpha\frac{1}{T}\sum_{t=1}^{T}
    D_{\mathrm{KL}}\left(
       p_T(\cdot\mid x_{<t})
       \Vert p_S(\cdot\mid x_{<t})
    \right).
    \label{eq:distillation-pretraining-objective}
\end{equation}
The central motivation is that the teacher distribution contains information
about alternatives that the observed next token does not reveal.  The
pre-training design-space study by \textcite{Peng2024PretrainingDistillation}
examines the processing of teacher logits, the loss function, the allocation
of distillation weight, model and data scaling, and offline versus online
generation of teacher targets.

\paragraph{Truncating teacher logits.}
A full teacher distribution has \(C\) entries at every token position.  If
there are \(N\) training tokens and each stored logit uses \(b\) bytes, the
storage required for all teacher logits is approximately
\begin{equation}
    M_{\mathrm{full}}=NCb.
    \label{eq:distillation-logit-storage}
\end{equation}
For billion-token corpora and vocabularies with hundreds of thousands of
entries, this can require petabytes of storage.  Peng et al. report that
storing full float32 logits for their 100-billion-token preliminary setting
would require approximately 58.6 petabytes.  Their two-stage truncation
procedure applies top-\(p\) selection followed by top-\(k\) selection and
renormalizes the retained probability mass.  In top-\(p\) selection, sort
the teacher probabilities in descending order and retain the smallest set
\(S_p\) satisfying
\begin{equation}
    \sum_{i\in S_p}p_{T,(i)}\geq p,
    \label{eq:distillation-top-p}
\end{equation}
where \(p_{T,(i)}\) denotes the \(i\)-th largest teacher probability.
Top-\(k\) then caps the number of retained entries at \(k\).  If the
retained index set is \(S\), the normalized truncated target is
\begin{equation}
    \widetilde p_{T,i}
    =\frac{p_{T,i}\mathbb I[i\in S]}
    {\displaystyle\sum_{j\in S}p_{T,j}},
    \label{eq:distillation-truncated-target}
\end{equation}
where \(\mathbb I[\cdot]\) is the indicator function.  Renormalization
is necessary because truncation removes probability mass.  The reported
top-0.95-100 configuration reduces storage by roughly a factor of 4,000,
from about 58.6 petabytes to approximately 15 terabytes in the preliminary
experiment.  Truncation therefore reduces the engineering burden, although
it also changes the target by removing the low-probability tail
\parencite{Peng2024PretrainingDistillation}.

\paragraph{Loss selection and scheduling.}
The pre-training study compares negative log-likelihood, KL divergence, and
mean squared error for distillation.  KL divergence and negative
log-likelihood perform similarly in the reported experiments, whereas mean
squared error is substantially less effective.  This distinction follows
from the object being matched: a probability distribution has a categorical
normalization constraint, while squared error measures Euclidean distance
between numerical coordinates and does not directly express a distributional
discrepancy.  The coefficient \(\alpha\) in
\cref{eq:distillation-pretraining-objective} also affects performance.
The reported experiments find that increasing the distillation contribution
helps up to a point, after which further increases can reduce performance.

Rather than holding \(\alpha\) constant, the study also examines a
warmup--stable--decay (WSD) schedule.  During warmup, \(\alpha\) increases
from zero to one; it remains at one during the stable phase and then follows
a cosine decay toward zero.  The reported configuration uses a warmup
proportion of approximately ten percent and a decay proportion of
approximately one percent, paired with a WSD learning-rate schedule.  The
schedule coordinates the strength of teacher supervision with the stage of
student optimization.  It is an empirical configuration, not a universal
schedule that must be optimal for every teacher, student, corpus, or training
budget \parencite{Peng2024PretrainingDistillation}.

\paragraph{Offline and online targets.}
Offline distillation computes teacher logits using an already pretrained
teacher and then uses those targets to train the student.  Online distillation
stores logits while the teacher itself is being pretrained.  Online storage
can avoid a separate teacher-inference pass and can make one teacher run
useful for a family of students.  However, the quality of online targets
depends on the teacher's training stage.  Peng et al. find that logits from
early, non-converged teacher checkpoints can be noisy and yield weak student
performance.  Later-stage online logits are more useful, although their
reported student remains below the student trained from offline logits of a
pretrained teacher.  Their practical conclusion is conditional: when one
student is required, offline targets from a pretrained teacher are preferable
in the studied setting; when multiple student sizes are to be trained, online
logits stored during the largest model's pre-training can amortize the target
generation cost \parencite{Peng2024PretrainingDistillation}.

The same study reports that larger students generally benefit more from
pre-training distillation, while increasing teacher size does not guarantee
better student performance.  Its preliminary experiment distills a
1.9-billion-parameter student from GLM-4-9B over 100 billion tokens and
reports an average relative improvement of approximately 1.6 percent over
the language-modeling baseline across the selected evaluations.  In its
larger-scale experiments, improvements remain visible over a 500-billion-
token training run.  These findings support the effectiveness of pre-training
distillation in the investigated settings, but do not establish a universal
gain at arbitrary model sizes or compute budgets
\parencite{Peng2024PretrainingDistillation}.

\subsection{On-policy distillation and the student-state distribution}
\label{subsec:distillation-on-policy}

The fixed-prefix objectives above define what the student should predict at
a given context, but they do not determine where that context comes from.
Autoregressive generation makes this distinction consequential.  Each
student-generated token becomes part of the next context, so an early
prediction changes the states at which later predictions are made.  A student
trained only on reference or teacher-generated prefixes may therefore be
poorly supervised on the histories it encounters during deployment.  On-policy
distillation (OPD) addresses this mismatch by generating student responses
and querying the teacher along the resulting student-generated prefixes
\parencite{Song2026OnPolicyDistillation}.

For a prompt \(x\) and a response \(y=(y_1,\ldots,y_T)\), define the
state at position \(t\) as
\begin{equation}
    c_t=(x,y_{<t}),
    \qquad y_{<t}=(y_1,\ldots,y_{t-1}).
    \label{eq:distillation-opd-state}
\end{equation}
The next token is the action.  The student policy is
\(\pi_\theta(a\mid c_t)\), and the teacher target is
\(q(a\mid c_t^T)\), where \(c_t^T\) denotes the context available to
the teacher.  In the simplest case, teacher and student are evaluated on
the same text prefix.  A trajectory consists of the sequence of states and
actions generated by the student.  Because each action changes the next
prefix, the state distribution is induced by the policy that generates the
response.

Under off-policy distillation, training states can come from reference
responses, teacher responses, or trajectories generated by an earlier
student.  Under strict on-policy distillation, the states are sampled from
the current student policy.  A trajectory collected by an earlier student
was on-policy when it was generated, but becomes stale relative to a later
student after the parameters change.  Re-scoring a stored trajectory can
refresh teacher supervision without refreshing the state distribution.
This distinction between the source of states and whether a dataset is stored
or processed offline is important when designing replay and asynchronous
training systems \parencite{Song2026OnPolicyDistillation}.

A common distribution-matching formulation is
\begin{equation}
    \mathcal L_{\mathrm{match}}(\theta)
    =\mathbb E_{s\sim d_\mu}
    \left[
      w(s)D_f\left(
      q(\cdot\mid s^T)\Vert
      \pi_\theta(\cdot\mid s)
      \right)
    \right],
    \label{eq:distillation-opd-general}
\end{equation}
where \(d_\mu\) is the state distribution induced by collection policy
\(\mu\), \(w(s)\) is a nonnegative state weight, and \(D_f\) is an
\(f\)-divergence.  For a convex generator \(f\) satisfying \(f(1)=0\),
\begin{equation}
    D_f(P\Vert Q)
    =\mathbb E_{a\sim Q}
    \left[
      f\left(\frac{P(a)}{Q(a)}\right)
    \right].
    \label{eq:distillation-f-divergence}
\end{equation}
Choosing \(f(u)=u\log u\) yields forward KL, while
\(f(u)=-\log u\) yields reverse KL, wherever the required probability
ratios are defined.  Equation \cref{eq:distillation-opd-general}
separates three decisions: which states are visited, which target the teacher
provides at those states, and which discrepancy is used to compare teacher
and student distributions.  Strict on-policy collection corresponds to
\(\mu=\pi_\theta\).

\paragraph{Sampled-token objectives and variance reduction.}
Full-vocabulary matching uses teacher probabilities for all possible next
tokens at a visited prefix.  A sampled-token surrogate can instead use only
the token drawn from the student.  Let
\begin{equation}
    r_t=\log\frac{p_T(y_t\mid c_t)}
    {p_\theta(y_t\mid c_t)}.
    \label{eq:distillation-opd-log-ratio}
\end{equation}
A detached reverse-KL surrogate is
\begin{equation}
    \ell_t(\theta)
    =-\operatorname{sg}[r_t]
    \log p_\theta(y_t\mid c_t),
    \label{eq:distillation-opd-sampled-surrogate}
\end{equation}
where \(\operatorname{sg}[\cdot]\) denotes stop-gradient: the
coefficient is held fixed while differentiating the student log-probability.
If \(r_t>0\), the teacher assigns the sampled token greater probability
than the student and the update reinforces that token.  If \(r_t<0\), the
update suppresses it.  This estimator avoids an explicit sum over the full
vocabulary, but it does not directly provide a positive correction for every
teacher-preferred alternative that the student did not sample.  Restricting
supervision to top-\(K\) tokens, discarding the tail, renormalizing it, or
representing it as a residual probability bucket changes the effective
support of the objective \parencite{Song2026OnPolicyDistillation}.

A baseline can reduce sampling variance without changing the expected
conditional score-function gradient when it depends on the prefix but not
on the sampled action.  The conditional mean log-ratio baseline is
\begin{align}
    b_t(c_t)
    &=\mathbb E_{a\sim p_\theta(\cdot\mid c_t)}
      \left[
      \log\frac{p_T(a\mid c_t)}{p_\theta(a\mid c_t)}
      \right]\nonumber\\
    &=-D_{\mathrm{KL}}\left(
      p_\theta(\cdot\mid c_t)
      \Vert p_T(\cdot\mid c_t)
      \right).
    \label{eq:distillation-opd-baseline}
\end{align}
Subtracting this action-independent baseline centers the immediate
conditional learning signal.  It is a control variate: it changes the
variance of the estimator without changing its expectation under the stated
sampling conditions.  It should not be confused with a value function that
estimates cumulative future discrepancies
\parencite{Song2026OnPolicyDistillation}.

\paragraph{Sequence-level reverse KL and temporal credit.}
A local conditional objective treats a visited prefix as fixed.  It does not
necessarily assign credit to an earlier token for changing the prefixes that
will be visited later.  Sequence-level reverse KL accounts for this temporal
dependence.  Let \(p_\theta(y)\) and \(p_T(y)\) be the student and teacher
sequence distributions, and let the student trajectory be sampled from
\(p_\theta\).  For a fixed teacher with support on the student-generated
responses, the ideal student-sampled gradient can be written as
\begin{equation}
    \nabla_\theta D_{\mathrm{KL}}(p_\theta\Vert p_T)
    =-\mathbb E_{y\sim p_\theta}
    \left[
      \sum_{t=1}^{|y|}(R_t-1)
      \nabla_\theta\log p_\theta(y_t\mid c_t)
    \right],
    \label{eq:distillation-sequence-reverse-kl-gradient}
\end{equation}
where the return-to-go of log-probability ratios is
\begin{equation}
    R_t=\sum_{t'=t}^{|y|}
    \log\frac{p_T(y_{t'}\mid c_{t'})}
    {p_\theta(y_{t'}\mid c_{t'})}.
    \label{eq:distillation-sequence-return}
\end{equation}
The constant score contribution vanishes in expectation, whereas \(R_t\)
allows an earlier action to receive credit for teacher--student discrepancies
at later positions.  MiniLLM uses sequence-level reverse KL and approximates
the practical update using teacher--student mixed sampling, per-token
importance ratios, and length normalization.  These implementation choices
matter: without suitable sampling and length control, reverse-KL training can
reduce its objective while producing short, repetitive, or otherwise
unhelpful responses \parencite{Song2026OnPolicyDistillation}.

A discounted temporal objective makes the horizon of credit assignment
explicit.  Define the local log-probability discrepancy
\begin{equation}
    \Delta_t
    =\log p_\theta(y_t\mid c_t)
    -\log p_T(y_t\mid c_t).
    \label{eq:distillation-opd-local-gap}
\end{equation}
The discounted return-like quantity is
\begin{equation}
    A_t^{(\gamma)}
    =-\sum_{t'=t}^{T}
    \gamma^{t'-t}\Delta_{t'},
    \qquad 0\leq\gamma\leq 1.
    \label{eq:distillation-opd-discounted-advantage}
\end{equation}
At \(\gamma=0\), the signal is local to the current token.  At
\(\gamma=1\), it includes the undiscounted future discrepancies through
the end of the sequence.  Intermediate values trade off local correction
against longer-range temporal credit.  Under the uniform second-moment
assumption analyzed in the on-policy survey, fixed \(\gamma<1\) gives a
horizon-independent upper bound on the variance of this advantage estimate;
that statement is not, by itself, a variance bound for the complete
parameter-gradient estimator \parencite{Song2026OnPolicyDistillation}.

The teacher's reliability at student-generated prefixes is a separate
condition.  A prefix may contain an incorrect intermediate step or a context
that the teacher would not ordinarily generate.  High standalone teacher
accuracy does not guarantee reliable guidance at every such state.  Teacher
adaptation, eligibility tests, target repair, and selective supervision are
therefore relevant extensions.  On-policy distillation solves an important
state-distribution mismatch, but it does not make every teacher prediction
correct or every teacher target learnable by the student
\parencite{Song2026OnPolicyDistillation}.

\begin{figure}[htbp]
    \centering
    % On-policy distillation loop diagram.
% Uses the standard LaTeX picture environment; no additional package is needed.
\begingroup
\setlength{\unitlength}{0.75mm}
\begin{picture}(150,66)
    \put(1,43){\framebox(25,12){Prompt}}
    \put(37,43){\framebox(30,12){\shortstack{Student\\rollout}}}
    \put(78,43){\framebox(30,12){\shortstack{Student\\prefixes}}}
    \put(119,43){\framebox(29,12){\shortstack{Teacher\\scoring}}}
    \put(119,12){\framebox(29,12){\shortstack{Teacher\\targets}}}
    \put(68,12){\framebox(34,12){\shortstack{Distillation\\update}}}
    \put(26,49){\vector(1,0){11}}
    \put(67,49){\vector(1,0){10}}
    \put(108,49){\vector(1,0){10}}
    \put(133,43){\vector(0,-1){18}}
    \put(119,18){\vector(-1,0){16}}
    \put(68,18){\vector(-1,1){18}}
    \put(37,35){\makebox(30,5)[c]{\scriptsize updated policy}}
    \put(2,3){\makebox(145,5)[l]{\scriptsize The updated student changes the distribution of prefixes in the next rollout.}}
\end{picture}
\endgroup

    \caption{The on-policy distillation loop.  The student generates its own
    response prefixes; the teacher evaluates those prefixes and supplies
    targets; the resulting update changes the student that generates the next
    trajectories.  The state distribution therefore changes as the student
    is optimized.}
    \label{fig:distillation-on-policy-loop}
\end{figure}

\subsection{Distillation scaling laws and compute allocation}
\label{subsec:distillation-scaling-laws}

A distillation objective specifies what the student should imitate, but it
does not guarantee that the student has enough representational capacity,
data, or compute to learn the target.  The teacher may be more capable than
the student can approximate.  Conversely, a teacher that is too weak may
provide little information beyond the original training targets.  The
teacher--student capacity gap must therefore be considered jointly with the
student's training budget.  The distillation scaling study by
\textcite{Busbridge2025DistillationScalingLaws} makes this relationship
quantitative through validation cross-entropy.

A conventional supervised scaling relation expresses validation loss as a
function of parameter count \(N\) and training tokens \(D\):
\begin{equation}
    L(N,D)
    =E+\left(\frac{A}{N^\alpha}
    +\frac{B}{D^\beta}\right)^\gamma,
    \label{eq:distillation-supervised-scaling-law}
\end{equation}
where \(E\) is an irreducible-error term under the fitted model,
\(A,B,\alpha,\beta,\gamma\) are positive fitted coefficients, and
\(N,D\) represent model capacity and data scale.  Increasing either model
size or training data generally reduces loss with diminishing marginal
returns.  Under a fixed compute budget, compute-optimal training chooses the
model size and token count that minimize validation loss subject to the
training FLOP constraint.

For distillation, let \(N_S\) be the number of student non-embedding
parameters, \(D_S\) the number of distillation tokens, and \(L_T\) the
teacher validation cross-entropy.  Let
\begin{equation}
    \widetilde L_S=L(N_S,D_S)
    \label{eq:distillation-supervised-student-reference}
\end{equation}
be the cross-entropy expected for a student of that size and token budget
under the supervised scaling relation.  Busbridge et al. fit the following
distillation scaling law:
\begin{equation}
\begin{aligned}
    L_S(N_S,D_S,L_T)
    ={}&L_T
    +c_0\left[
      1+\left(\frac{L_T}{\widetilde L_S}\right)^{1/f_1}
    \right]^{-c_1f_1}\\
    &\times\left(
      \frac{A}{N_S^{\alpha'}}
      +\frac{B}{D_S^{\beta'}}
    \right)^{\gamma'},
\end{aligned}
    \label{eq:distillation-scaling-law}
\end{equation}
where \(c_0,c_1,f_1,A,B,\alpha',\beta',\gamma'\) are positive
coefficients fitted to the experimental data.  The first term is the
teacher cross-entropy.  The second term models the student's ability to
approximate the teacher and depends on student size and distillation data.
The ratio \(L_T/\widetilde L_S\) captures the relative learning abilities
of teacher and student on the reference task.  The fitted transition
produces a capacity gap: at fixed student size and data budget, improving
the teacher initially helps, but a teacher that becomes too difficult for
the student to imitate can eventually produce a worse student.  The optimal
teacher is therefore not necessarily the strongest available teacher.

The scaling law is empirical and was fitted to a particular family of
language models and training conditions.  It should not be interpreted as a
universal law independent of architecture, data distribution, tokenization,
or optimization.  Its main contribution is to make teacher quality,
student capacity, and the number of distillation tokens explicit in one
predictive relation \parencite{Busbridge2025DistillationScalingLaws}.

\paragraph{Compute accounting.}
Distillation consumes compute for student training, teacher inference, and,
when the teacher does not already exist, teacher pre-training.  Let \(F(N)\)
be the forward-pass FLOPs per token for a model with \(N\) parameters.
Busbridge et al. approximate the total cost as
\begin{equation}
    C_{\mathrm{distill}}
    \approx 3F(N_S)D_S
    +F(N_T)\left(
       \delta_{T}^{\mathrm{logit}}D_S
       +3\delta_{T}^{\mathrm{pre}}D_T
    \right),
    \label{eq:distillation-compute-cost}
\end{equation}
where \(N_T\) and \(D_T\) are teacher parameter count and pre-training
tokens.  The factor of three in the student term approximates forward and
backward training cost; the teacher-logit term accounts for teacher forward
passes used to produce targets; and the teacher-pretraining term accounts
for training the teacher.  The indicators
\(\delta_T^{\mathrm{logit}}\) and
\(\delta_T^{\mathrm{pre}}\) specify which costs are charged to the
current budget.  A fully amortized teacher can make both indicators zero;
if the teacher already exists but must be queried for this student, only
the logit cost is counted; if teacher training is required and its logits
are stored for reuse, the pre-training cost is counted; if both teacher
training and student-specific inference are charged, both terms are
included.

This accounting changes the comparison with supervised learning.  When a
suitable teacher already exists or will be reused across many students,
distillation can amortize the teacher cost and outperform supervised
training over a range of budgets.  When a teacher must be trained solely
for one student, the teacher's pre-training and inference costs can make
supervised learning preferable.  Busbridge et al. find that the compute
regime in which distillation is advantageous depends on student size and
teacher reuse; distillation is not automatically more efficient merely
because the final student is smaller
\parencite{Busbridge2025DistillationScalingLaws}.

\subsection{Dataset distillation and its integration with model distillation}
\label{subsec:distillation-dataset-distillation}

Knowledge distillation compresses a model's learned behavior.  Dataset
distillation addresses a complementary problem: it attempts to construct a
small synthetic dataset that preserves information needed to train a model.
Let the original dataset be
\(D_R=\{(x_i,y_i)\}_{i=1}^{N}\), and let a compact synthetic dataset be
\(D_S=\{(\widetilde x_j,\widetilde y_j)\}_{j=1}^{M}\), with
\(M\ll N\).  Unlike ordinary data selection, which retains existing
examples, dataset distillation can optimize or generate synthetic examples
specifically to preserve the learning signal of the original corpus
\parencite{Fang2026KnowledgeDatasetDistillation}.

\paragraph{Gradient-matching formulations.}
Let \(\mathcal L(D;\theta)\) be the training loss of a model with
parameters \(\theta\) on dataset \(D\).  At optimization state
\(\theta_t\), define the original-data and synthetic-data gradients as
\begin{equation}
    g_R^{(t)}=\nabla_{\theta_t}\mathcal L(D_R;\theta_t),
    \qquad
    g_S^{(t)}=\nabla_{\theta_t}\mathcal L(D_S;\theta_t).
    \label{eq:distillation-dataset-gradients}
\end{equation}
A gradient-matching objective seeks a synthetic dataset whose updates
resemble those induced by the original data:
\begin{equation}
    \min_{D_S}
    \sum_{t=0}^{K-1}
    \left\|
      \nabla_{\theta_t}\mathcal L(D_S;\theta_t)
      -\nabla_{\theta_t}\mathcal L(D_R;\theta_t)
    \right\|_2^2.
    \label{eq:distillation-gradient-matching}
\end{equation}
The Euclidean norm measures the discrepancy between the parameter updates
induced by the two datasets.  If the synthetic examples produce similar
gradients over the relevant training trajectory, training on the compact
set may approximate aspects of training on the full dataset.  The method
can be computationally expensive because synthetic samples must be
optimized against model-training dynamics, often involving repeated
higher-order or unrolled computations.

\paragraph{Generative synthesis and selection.}
A generative approach learns a mapping \(G\) from a latent variable
\(z\sim p(z)\) to a synthetic example:
\begin{equation}
    \widetilde x=G(z),
    \qquad z\sim p(z).
    \label{eq:distillation-generative-sample}
\end{equation}
The generator is trained to produce representative examples of the original
data distribution or examples that preserve task-relevant information.  The
survey also distinguishes dataset synthesis from data selection: selection
identifies a high-quality subset of existing examples, whereas synthesis
constructs new examples through optimization, generative models, or
heuristics.  These methods trade off computational cost, sample diversity,
and the fidelity with which rare or difficult patterns are preserved
\parencite{Fang2026KnowledgeDatasetDistillation}.

\paragraph{Combining dataset and knowledge distillation.}
The two forms of compression can be integrated.  A teacher is first trained
on the original dataset; a compact synthetic coreset is then constructed;
and teacher predictions can be used as soft targets on the synthetic
examples.  If
\(C_{\mathrm{syn}}=\{(\widetilde x_i,\widetilde y_i)\}_{i=1}^{M}\)
is the synthetic coreset, a student can be trained using
\begin{equation}
    \mathcal L_{C_{\mathrm{syn}}}(\theta_S)
    =-\frac{1}{M}\sum_{i=1}^{M}
    \sum_{v\in\mathcal V}
    \widetilde y_{i,v}
    \log p_{\theta_S}(v\mid\widetilde x_i),
    \label{eq:distillation-synthetic-soft-label-loss}
\end{equation}
where \(\widetilde y_i\) is the teacher-derived target distribution.
This combines compression of the training examples with compression of the
teacher's predictive behavior.  The central risk is that aggressive dataset
compression can remove rare linguistic structures, facts, or reasoning
patterns that occupy little volume in the original corpus but matter for
generalization.  The size of the synthetic set is therefore not sufficient
evidence of its information coverage
\parencite{Fang2026KnowledgeDatasetDistillation}.

\subsection{Evaluation of distilled language models}
\label{subsec:distillation-evaluation}

A distilled model must be evaluated along several dimensions because
predictive quality, teacher fidelity, efficiency, and robustness are not
equivalent objectives.  A student can reproduce teacher outputs closely
without generalizing well to unseen tasks.  It can also outperform the
teacher on a particular benchmark while differing substantially from the
teacher's original distribution.  Evaluation must therefore compare the
student both with the teacher and with independent task targets
\parencite{Yang2024SurveyKD,Fang2026KnowledgeDatasetDistillation}.

For language modeling, perplexity is the exponential of token-averaged
negative log-likelihood:
\begin{equation}
    \operatorname{PPL}
    =\exp\left(
      -\frac{1}{T}\sum_{t=1}^{T}
      \log p_{\theta_S}(x_t\mid x_{<t})
    \right).
    \label{eq:distillation-perplexity}
\end{equation}
Lower perplexity means that the model assigns higher average probability
to the observed evaluation tokens.  For task-based evaluation, accuracy,
exact match, and F1-score measure different aspects of answer correctness.
Benchmarks for knowledge, reasoning, mathematics, commonsense, and
instruction following are needed to determine whether the student preserves
capabilities beyond the examples used for distillation.  A benchmark result
is meaningful only with its task definition, data split, prompting protocol,
and scoring procedure specified.

Teacher fidelity can be measured directly through a distributional
comparison such as
\begin{equation}
    \mathcal F_{\mathrm{KD}}
    =\mathbb E_{x\sim\mathcal D_{\mathrm{eval}}}
    \left[
      D_{\mathrm{KL}}\left(
      p_T(\cdot\mid x)\Vert p_S(\cdot\mid x)
      \right)
    \right].
    \label{eq:distillation-fidelity}
\end{equation}
A lower value indicates closer predictive distributions under the selected
evaluation contexts.  Representation-based methods can additionally compare
hidden states using cosine similarity or representation-similarity measures.
These metrics assess transfer fidelity, not necessarily task correctness.

Efficiency evaluation should include model parameter count, storage,
peak memory, inference latency, throughput, and arithmetic cost.  A simple
parameter compression ratio is
\begin{equation}
    R_{\mathrm{compression}}
    =\frac{N_T}{N_S},
    \label{eq:distillation-compression-ratio}
\end{equation}
but parameter count alone does not determine wall-clock latency or memory
consumption.  Architecture, sequence length, hardware, batching, and
implementation all influence actual deployment cost.  Training efficiency
must also account for teacher inference and, where relevant, teacher
pre-training, as in \cref{eq:distillation-compute-cost}.

Robustness and calibration are also relevant when the student is expected
to operate under noisy or shifted inputs.  Predictive entropy
\cref{eq:distillation-entropy} describes the uncertainty of a single
categorical distribution, while calibration compares stated confidence
with empirical accuracy.  One common calibration statistic is expected
calibration error (ECE).  If evaluation examples are partitioned into
confidence bins \(B_m\), then
\begin{equation}
    \operatorname{ECE}
    =\sum_{m=1}^{M}
    \frac{|B_m|}{N}
    \left|
      \operatorname{acc}(B_m)
      -\operatorname{conf}(B_m)
    \right|.
    \label{eq:distillation-ece}
\end{equation}
Here \(\operatorname{acc}(B_m)\) is empirical accuracy in bin \(m\),
and \(\operatorname{conf}(B_m)\) is its average predicted confidence.
Robustness testing should also consider distribution shifts and adversarial
perturbations where appropriate.  No single metric captures all these
properties; a distillation method should be assessed against the deployment
requirements it is intended to satisfy
\parencite{Yang2024SurveyKD,Fang2026KnowledgeDatasetDistillation}.

\subsection{Limitations and connection to parameter-efficient fine-tuning}
\label{subsec:distillation-limitations-next-topic}

Knowledge distillation reduces deployment cost only when the student
preserves the capabilities needed by its application.  Several limitations
follow from the mechanisms developed above.  A narrow distillation dataset
may omit capabilities that are not expressed in its prompts or examples.
Teacher-generated rationales may be fluent but incorrect, and a student can
inherit teacher biases or systematic errors.  Multi-teacher training can
combine complementary expertise but also increases inference cost and
creates conflicts among heterogeneous supervisors.  Hint-based methods
require decisions about which layers or representations should be aligned,
and intermediate activation storage can become expensive.  On-policy
methods expose the teacher to student-generated histories, but teacher
reliability at erroneous or unfamiliar prefixes remains a separate problem.
Finally, a teacher that is substantially more capable than the student may
be difficult to imitate, while the cost of producing teacher targets can
outweigh the savings from a smaller student
\parencite{Yang2024SurveyKD,Peng2024PretrainingDistillation,Busbridge2025DistillationScalingLaws,Song2026OnPolicyDistillation}.

These limitations show why distillation is not a universal replacement for
conventional pre-training or fine-tuning.  Its value depends on the
information available from the teacher, the student's capacity, the data
coverage, the optimization objective, and the total compute budget.  Dataset
distillation can reduce the volume of training information that must be
processed, while quantization reduces the numerical precision required for
model storage and computation.  Distillation addresses a different source
of cost: it transfers learned behavior to a model with a smaller or more
suitable parameterization.  These strategies can be combined, but their
benefits must be measured at the level of the complete training and
inference pipeline.

The next section turns from transferring behavior between models to
adapting an existing model with a restricted set of trainable parameters.
Parameter-efficient fine-tuning investigates how a pretrained model can be
specialized without updating all of its parameters.  It therefore continues
the chapter's central optimization question from a complementary direction:
how can useful model behavior be obtained and adapted while controlling the
computational and memory costs of learning?
